\section{Sliding Window Mex}
\label{sec:cses-sw-mex}

\problemstatement{Given an array of $n$ non-negative integers and a window
size $k$, output the mex of every contiguous window -- the smallest
non-negative integer \emph{not} present in the window.}

\begin{patternbox}
The mex of a size-$k$ window is at most $k$ (with $k$ elements, some value
in $0\ldots k$ must be missing). So track only values in $[0,k]$ with a
frequency map and keep an ordered \textbf{set of absent values}; the mex is
its smallest element.
\end{patternbox}

\subsection*{A set of the missing small values}

Values larger than $k$ can never be the mex, so ignore them. Initialise a
set $\textit{absent}=\{0,1,\ldots,k\}$ and a frequency map. When a value
$v\le k$ enters and its frequency goes $0\to1$, erase $v$ from
$\textit{absent}$ (now present). When it leaves and its frequency returns
to $0$, insert $v$ back into $\textit{absent}$. The mex is the smallest
element of $\textit{absent}$ -- read in $O(\log k)$ from its begin.

\begin{ideabox}[Bound the Universe, Then Track Absence]
Because the mex cannot exceed the window size, the relevant universe is
just $[0,k]$. Maintaining which of those are absent -- rather than which are
present -- puts the answer (the smallest absent value) at the front of an
ordered set.
\end{ideabox}

\subsection*{Algorithm}

\begin{enumerate}
  \item $\textit{absent}=\{0,\ldots,k\}$; empty frequency map.
  \item Add the first $k$ elements (erasing from $\textit{absent}$ on
        $0\to1$).
  \item Each slide: add $a[i]$, remove $a[i-k]$ (inserting into
        $\textit{absent}$ on $1\to0$); output $\ast\textit{absent}.
        \text{begin}()$.
\end{enumerate}

\subsection*{Dry run}

\dryrun{$a=[0,1,3,0]$, $k=3$.}

\begin{itemize}
  \item $\textit{absent}=\{0,1,2,3\}$. Window $[0,1,3]$: erase $0,1,3$;
        absent $=\{2\}$; mex $2$.
  \item Slide to $[1,3,0]$: remove first $0$ (freq $0$, reinsert $0$),
        add $0$ (erase $0$ again); absent $=\{2\}$; mex $2$.
  \item Output $2,2$.
\end{itemize}

\subsection*{C++ implementation}

\begin{lstlisting}
#include <bits/stdc++.h>
using namespace std;

int main() {
    int n, k;
    cin >> n >> k;
    vector<int> a(n);
    for (int& x : a) cin >> x;

    set<int> absent;
    for (int v = 0; v <= k; v++) absent.insert(v);
    unordered_map<int,int> freq;

    auto add = [&](int v) {
        if (v <= k && freq[v]++ == 0) absent.erase(v);
    };
    auto remove = [&](int v) {
        if (v <= k && --freq[v] == 0) absent.insert(v);
    };

    for (int i = 0; i < k; i++) add(a[i]);
    cout << *absent.begin() << " ";
    for (int i = k; i < n; i++) {
        add(a[i]); remove(a[i - k]);
        cout << *absent.begin() << " ";
    }
    cout << "\n";
    return 0;
}
\end{lstlisting}

\begin{complexitybox}
\textbf{Time Complexity.} $O(n\log k)$ -- ordered-set operations bounded by
the $[0,k]$ universe. \textbf{Space Complexity.} $O(k)$.
\end{complexitybox}

\begin{takeawaybox}
Mex problems shrink to a bounded universe (here $[0,k]$) and then track
\emph{absence} in an ordered set so the answer is the set's minimum. Bounding
the universe before choosing a structure is the recurring first move for
mex.
\end{takeawaybox}
